% ══════════════════════════════════════════════════════════════════
%  FICHE D'EXERCICES — Chapitre 11 : Acides et bases — pH, couples, pKa
% ══════════════════════════════════════════════════════════════════
\section{Exercices type Baccalauréat}
\textit{Données à \SI{25}{\celsius} : $pK_e = 14$ ; $pH = -\log[\ce{H3O+}]$. On néglige
les ions provenant de l'eau sauf indication contraire.}

\rubrique{Couples acide/base, équations}

\exo{}\diff{1} Donner le couple acide/base associé à chacune des espèces :
a)~$\ce{NH3}$ ; b)~$\ce{CH3COOH}$ ; c)~$\ce{HCO3-}$ ; d)~$\ce{H2O}$.

\exo{}\diff{1} Écrire l'équation de la réaction entre l'acide éthanoïque et
l'ammoniac. Identifier les deux couples acide/base mis en jeu.

\exo{}\diff{1} Écrire l'équation d'autoprotolyse de l'eau et donner l'expression du
produit ionique $K_e$.

\rubrique{Calculs de pH — acides et bases forts}

\exo{}\diff{1} Calculer le pH d'une solution d'acide chlorhydrique de concentration
$C = \SI{e-2}{\mol\per\litre}$, puis celui d'une solution d'hydroxyde de sodium de
même concentration.

\exo{}\diff{2} On dissout \SI{0.40}{\gram} d'hydroxyde de sodium ($M=\SI{40}{\gram\per\mol}$)
dans l'eau pour obtenir \SI{500}{\milli\litre} de solution.
\begin{enumerate}[label=\arabic*)]
\item Calculer la concentration de la solution.
\item En déduire $[\ce{HO-}]$, $[\ce{H3O+}]$ et le pH.
\end{enumerate}

\exo{}\diff{2} À \SI{100}{\milli\litre} d'une solution d'acide chlorhydrique de pH $=2$,
on ajoute \SI{900}{\milli\litre} d'eau. Calculer le nouveau pH. Conclure sur l'effet
d'une dilution au dixième sur le pH d'un acide fort.

\rubrique{Acides et bases faibles, $K_a$ et $pK_a$}

\exo{}\diff{2} Une solution d'acide éthanoïque de concentration
$C = \SI{e-2}{\mol\per\litre}$ a un pH de $3{,}4$.
\begin{enumerate}[label=\arabic*)]
\item Montrer que l'acide éthanoïque est un acide faible (calculer le coefficient de
dissociation $\alpha$).
\item Déterminer les concentrations de toutes les espèces.
\item Calculer le $K_a$ et le $pK_a$ du couple $\ce{CH3COOH}/\ce{CH3COO-}$.
\end{enumerate}

\exo{}\diff{2} Le couple $\ce{NH4+}/\ce{NH3}$ a pour $pK_a = 9{,}2$. Tracer le
diagramme de prédominance et indiquer l'espèce prédominante à pH $= 6$, à pH $= 9{,}2$
et à pH $= 11$.

\exo{}\diff{3} Une solution d'acide méthanoïque (formique) $\ce{HCOOH}$ de
concentration $C = \SI{0.10}{\mol\per\litre}$ a un pH de $2{,}4$.
\begin{enumerate}[label=\arabic*)]
\item Faire le bilan des espèces et écrire les équations de conservation et
d'électroneutralité.
\item Déterminer les concentrations des espèces présentes.
\item En déduire le $K_a$ et le $pK_a$ du couple, puis comparer à celui de l'acide
éthanoïque ($pK_a = 4{,}8$). Lequel est le plus fort ?
\end{enumerate}

\rubrique{Problème de synthèse — type Baccalauréat}

\sujetbac[d'après Baccalauréat C, Cameroun]
On dispose d'une solution $S$ d'un monoacide $AH$ de concentration
$C = \SI{2.0e-2}{\mol\per\litre}$ et de pH $= 3{,}0$.
\begin{enumerate}[label=\arabic*)]
\item Calculer $[\ce{H3O+}]$ et $[\ce{HO-}]$ dans la solution.
\item L'acide $AH$ est-il fort ou faible ? Justifier par le calcul.
\item Faire le bilan des espèces et déterminer toutes les concentrations.
\item Déterminer le $pK_a$ du couple $AH/A^-$.
\item On ajoute de l'eau jusqu'à diluer $S$ dix fois. Le pH augmente-t-il de une
unité ? Justifier qualitativement.
\end{enumerate}
